\makeatletter

% --- MISE EN PAGES --- %
\usepackage[T1]{fontenc}
\usepackage[utf8]{inputenc}
\usepackage[french]{babel}
\usepackage{xcolor}
\usepackage{lmodern}
\usepackage{dsfont}
\usepackage[a4paper,hmargin=20mm,vmargin=20mm]{geometry}

\usepackage{subcaption}
\usepackage{tocloft}

\usepackage{yhmath}
\usepackage{eurosym}
\usepackage{diagbox}
\usepackage{enumitem}
\setlist[itemize,1]{label={\color{black}\small \textbullet}}
\usepackage{multicol}

\usepackage{graphicx}
\usepackage{float}
\usepackage{minted}
\usepackage{tikz}

\usepackage{array}
\newcolumntype{L}[1]{>{\raggedright\let\newline\\\arraybackslash\hspace{0pt}}m{#1}}
\newcolumntype{C}[1]{>{\centering\let\newline\\\arraybackslash\hspace{0pt}}m{#1}}
\newcolumntype{R}[1]{>{\raggedleft\let\newline\\\arraybackslash\hspace{0pt}}m{#1}}


% --- DEFINE COLORS --- %
\definecolor{astral}{RGB}{46,116,181}

\definecolor{lightorange}{RGB}{255,176,85} % couleur définition poly Dupont
\definecolor{lightblue}{RGB}{137,181,226} % couleur proposition poly Dupont
\definecolor{dustyrose}{RGB}{191,128,128} % couleur exemple poly Dupont
\definecolor{midblue}{RGB}{42,112,180} % couleur théorème poly Dupont
\definecolor{lightpink}{RGB}{255,170,170} % couleur exercice poly Dupont

\definecolor{customlightblue}{RGB}{145,200,235} % bleu plus clair que celui des propositions
\definecolor{matchinggreen}{RGB}{120,190,150} % vert dans le même ton que le bleu


%----------
%  Version
%----------
\usepackage[colorlinks=true, linkcolor=astral]{hyperref}
\hypersetup{linktocpage, colorlinks=true, linkcolor=astral}

\usepackage{fancyhdr} 
\usepackage{stmaryrd}


%----------
% Maths environnement 
%----------
\usepackage{latexsym}
\usepackage{amsmath}
\usepackage{amsbsy}
\usepackage{amsfonts}
\usepackage{amssymb}
\usepackage{nicefrac}
\usepackage{amscd}
\usepackage{amsthm}
\usepackage{mathtools}
\usepackage{tikz-cd}



% Commande norme triple
\newcommand{\vertiii}[1]{{\left\vert\kern-0.25ex\left\vert
    \kern-0.25ex\left\vert #1 
\right\vert\kern-0.25ex\right\vert\kern-0.25ex\right\vert}}

% Commande quotient propre
\newcommand{\quotient}[2]{#1/{#2}}

% Commande exponentielle
\newcommand{\ex}{\mathrm{e}}

% d droit
\renewcommand{\d}{\mathrm{d}}

% L^1_{loc}
\newcommand{\Lloc}[1]{L^1_{\text{loc}}\big(#1\big)}

% partie réelle
\renewcommand{\Re}{\mathfrak{Re}}

% partie imaginaire
\renewcommand{\Im}{\mathfrak{Im}}

% image
\newcommand{\im}{\mathrm{Im}}

\numberwithin{equation}{section}

\newtheoremstyle{definitionSs}{\topsep}{\topsep}%
     {}%         Body font
     {}%         Indent amount (empty = no indent, \parindent = para indent)
     {\sffamily\bfseries}% Thm head font
     {.}%        Punctuation after thm head
     { }%     Space after thm head (\newline = linebreak)
     {\thmname{#1}\thmnumber{~#2}\thmnote{~#3}}%         Thm head spec

\newtheoremstyle{plainSs}{\topsep}{\topsep}%
     {\itshape}%         Body font
     {}%         Indent amount (empty = no indent, \parindent = para indent)
     {\sffamily\bfseries}% Thm head font
     {.}%        Punctuation after thm head
     { }%     Space after thm head (\newline = linebreak)
     {\thmname{#1}\thmnumber{~#2}\thmnote{~#3}}%         Thm head spec

\theoremstyle{definitionSs}
\newtheorem{remark}{Remarque}[section]
\newtheorem{example}{Exemple}[section]

%\usepackage[framemethod=tikz]{mdframed}
\usepackage[]{mdframed}
\usepackage[nameinlink,noabbrev]{cleveref}

\newmdtheoremenv[
hidealllines=true,
leftline=true,
skipabove=0pt,
innertopmargin=-5pt,
innerbottommargin=2pt,
linewidth=4pt,
linecolor=lightorange,
innerrightmargin=0pt,
]{definition}{Définition}[section]

\newmdtheoremenv[
hidealllines=true,
leftline=true,
skipabove=0pt,
innertopmargin=-5pt,
innerbottommargin=2pt,
linewidth=4pt,
linecolor=dustyrose,
innerrightmargin=0pt,
]{propdef}[definition]{Définition - Proposition}

\newmdtheoremenv[
hidealllines=true,
leftline=true,
skipabove=0pt,
innertopmargin=-5pt,
innerbottommargin=2pt,
linewidth=4pt,
linecolor=midblue,
innerrightmargin=0pt,
]{theorem}[definition]{Théorème}

\newmdtheoremenv[
hidealllines=true,
leftline=true,
skipabove=0pt,
innertopmargin=-5pt,
innerbottommargin=2pt,
linewidth=4pt,
linecolor=matchinggreen,
innerrightmargin=0pt,
]{lemme}[definition]{Lemme}

\newmdtheoremenv[
hidealllines=true,
leftline=true,
skipabove=0pt,
innertopmargin=-5pt,
innerbottommargin=2pt,
linewidth=4pt,
linecolor=lightblue,
innerrightmargin=0pt,
]{proposition}[definition]{Proposition}

\newmdtheoremenv[
hidealllines=true,
leftline=true,
skipabove=0pt,
innertopmargin=-5pt,
innerbottommargin=2pt,
linewidth=4pt,
linecolor=customlightblue,
innerrightmargin=0pt,
]{corollaire}[definition]{Corollaire}

\newmdtheoremenv[
hidealllines=true,
leftline=true,
skipabove=0pt,
innertopmargin=-5pt,
innerbottommargin=2pt,
linewidth=4pt,
linecolor=lightpink,
innerrightmargin=0pt,
]{method}[definition]{Méthode}

% patch the way counter behave : one counter for all

\crefname{definition}{defi}{defis}

\crefname{propdef}{propdefs}{propdefss}
\let\oldpropdef\propdef
\renewcommand{\propdef}{%
  \crefalias{definition}{propdef}% definition counter now looks like lemme
  \oldpropdef}

\crefname{lemme}{lemma}{lemmas}
\let\oldlemme\lemme
\renewcommand{\lemme}{%
  \crefalias{definition}{lemme}% definition counter now looks like lemme
  \oldlemme}

\crefname{theorem}{thmm}{thmms}
\let\oldtheorem\theorem
\renewcommand{\theorem}{%
  \crefalias{definition}{theorem}% definition counter now looks like lemme
  \oldtheorem}

\crefname{proposition}{prop}{props}
\let\oldproposition\proposition
\renewcommand{\proposition}{%
  \crefalias{definition}{proposition}% definition counter now looks like lemme
  \oldproposition}

\crefname{corollaire}{corss}{corssss}
\let\oldcorollaire\corollaire
\renewcommand{\corollaire}{%
  \crefalias{definition}{corollaire}% definition counter now looks like lemme
  \oldcorollaire}


%---------------
% Mise en page
%--------------

\setlength{\parindent}{0pt}

\renewcommand*{\descriptionlabel}[1]{\hspace\labelsep{\sffamily #1}}
\providecommand{\defemph}[1]{{\sffamily\bfseries\color{black}#1}}
\renewcommand{\emph}[1]{{\sffamily #1}}

%\usepackage{titlesec, blindtext, color}
%\newcommand{\hsp}{\hspace{20pt}}
%\titleformat{\section}[hang]{\sffamily\Huge\bfseries}{\thesection\hsp\textcolor{gray75}{|}\hsp}{0pt}{\Huge\bfseries}
%\usepackage{sectsty}
%\allsectionsfont{\color{black}\normalfont\sffamily\bfseries}
%\renewcommand{\cfttoctitlefont}{\sffamily\bfseries\Huge\color{black}}
%\renewcommand{\cftchapfont}{\sffamily\bfseries\color{black}}
%\renewcommand{\cftsecfont}{\sffamily\bfseries}
%\renewcommand{\cftsubsecfont}{\sffamily}

% --- BIBLIOGRAPHIE --- %
\usepackage[backend=biber, style=numeric]{biblatex}
\addbibresource{mise_en_page/bibliographie/stageM1.bib}

\usepackage{csquotes}
\usepackage[toc, page]{appendix}


% -------------
% Some commands
% -------------

% --- LETTRES --- %
% --- \mathbb{}
\newcommand{\bbC}{\mathbb{C}}
\newcommand{\bbE}{\mathbb{E}}
\newcommand{\bbF}{\mathbb{F}}
\newcommand{\bbG}{\mathbb{G}}
\newcommand{\bbK}{\mathbb{K}}
\newcommand{\bbL}{\mathbb{L}}
\newcommand{\bbN}{\mathbb{N}}
\newcommand{\bbP}{\mathbb{P}}
\newcommand{\bbQ}{\mathbb{Q}}
\newcommand{\bbR}{\mathbb{R}}
\newcommand{\bbS}{\mathbb{S}}
\newcommand{\bbT}{\mathbb{T}}
\newcommand{\bbU}{\mathbb{U}}
\newcommand{\bbZ}{\mathbb{Z}}

% --- \mathbf{}
\newcommand{\bfC}{\mathbf{C}}
\newcommand{\bfE}{\mathbf{E}}
\newcommand{\bfL}{\mathbf{L}}
\newcommand{\bfN}{\mathbf{N}}
\newcommand{\bfP}{\mathbf{P}}
\newcommand{\bfQ}{\mathbf{Q}}
\newcommand{\bfR}{\mathbf{R}}
\newcommand{\bfZ}{\mathbf{Z}}

% --- \mathcal{}
\newcommand{\calA}{\mathcal{A}}
\newcommand{\calB}{\mathcal{B}}
\newcommand{\calC}{\mathcal{C}}
\newcommand{\calD}{\mathcal{D}}
\newcommand{\calE}{\mathcal{E}}
\newcommand{\calF}{\mathcal{F}}
\newcommand{\calG}{\mathcal{G}}
\newcommand{\calI}{\mathcal{I}}
\newcommand{\calJ}{\mathcal{J}}
\newcommand{\calL}{\mathcal{L}}
\newcommand{\calM}{\mathcal{M}}
\newcommand{\calN}{\mathcal{N}}
\newcommand{\calO}{\mathcal{O}}
\newcommand{\calP}{\mathcal{P}}
\newcommand{\calS}{\mathcal{S}}
\newcommand{\calT}{\mathcal{T}}
\newcommand{\calV}{\mathcal{V}}
\newcommand{\calX}{\mathcal{X}}
\newcommand{\calU}{\mathcal{U}}

% --- \mathfrak{}
\newcommand{\fA}{\mathfrak{A}}
\newcommand{\fB}{\mathfrak{B}}
\newcommand{\fF}{\mathfrak{F}}

% --- \mathscr{}
\usepackage{mathrsfs}
\newcommand{\scrL}{\mathscr{L}}
\newcommand{\scrF}{\mathscr{F}}
\newcommand{\scrM}{\mathscr{M}}
\newcommand{\scrD}{\mathscr{D}}
\newcommand{\scrE}{\mathscr{E}}
\newcommand{\scrA}{\mathscr{A}}
\newcommand{\scrB}{\mathscr{B}}
\newcommand{\scrC}{\mathscr{C}}
\newcommand{\scrI}{\mathscr{I}}
\newcommand{\scrJ}{\mathscr{J}}
\newcommand{\scrK}{\mathscr{K}}
\newcommand{\scrS}{\mathscr{S}}
\newcommand{\scrO}{\mathscr{O}}
\newcommand{\scrP}{\mathscr{P}}
\newcommand{\scrR}{\mathscr{R}}
\newcommand{\scrT}{\mathscr{T}}
\newcommand{\scrZ}{\mathscr{Z}}





